% ================================================================
\chapter{How to Play}
% ================================================================
% REAL RULES (from Notion + author): the core dice engine. Keep the
% mechanics verbatim; surrounding prose is open to editing.
\begin{readaloud}
\lipsum[5][1-3]
\end{readaloud}

\lettrine{L}{egend of the Technomancers} runs on one simple roll. Whenever
the outcome of an action is in doubt --- and only then --- you gather a
handful of six-sided dice, roll them all at once, and count your
\keyword{successes}. Everything else in this book is just a way of deciding
how many dice you get to throw.

\artbanner{play-dice.png}

\section{The Core Roll}
When a Traveller acts, say what they want and what they do. If the
outcome is certain, it happens; if it is impossible, the GM explains
why. When the outcome is in doubt and something is at stake, the GM sets
a target number and the player builds a pool of six-sided dice.

\begin{statblock}[Build Your Dice Pool]
\begin{center}
{\large\bfseries Base Stat + applicable tags = dice pool}
\end{center}
Roll the pool. Each die showing \keyword{4, 5, or 6} is one success.
Fewer successes than the target is a failure; exactly the target is a
partial success; more is a full success. If the number of \dice{6}s
alone exceeds the target, it is a critical success instead. See
\emph{Degrees of Success} in \emph{The Shape of a Session} for what
each result means.
\end{statblock}

\begin{example}
Mara jumps a gap. Her \keyword{Body} is \dice{3}, and her
\atrtag{skill}{Athletics} tag adds \dice{2} dice (one for the tag, one
for being her first Skill): she rolls \dice{5} dice. The GM set a target of \dice{2}. Mara gets three successes, so
she lands cleanly; with exactly two, she would land with a complication.
\end{example}

\section{Base Stat}
% REAL RULE (Notion, 2026-09-21, decided with Chris): THREE broad base
% stats — Mind, Body, Face. No separate Magic/Arcane stat: magic hangs off
% these same three, which stops one versatile stat from dominating play.
Every character has three broad base \keyword{Stats}. Everyone begins with
\dice{1} in each; character creation can raise them as high as \dice{5}.
Choose the one that fits what pushes back against the action, not just
how the action looks. A spell and a mundane action use the same Stats.

\begin{statblock}[The Three Stats]
\begin{tabularx}{\linewidth}{@{}p{0.30\linewidth} Y@{}}
\leavevmode{\color{\realmaccent}\atricon{body}}\ \keyword{Body} & Physical action and the handling of objects --- striking, lifting, running, throwing.\\[3pt]
\leavevmode{\color{\realmaccent}\atricon{mind}}\ \keyword{Mind} & Puzzles, analysis, knowledge, and perception.\\[3pt]
\leavevmode{\color{\realmaccent}\atricon{face}}\ \keyword{Face} & Dealing with people --- persuasion, deception, presence, command.\\
\end{tabularx}
\end{statblock}

On the page and on your sheet, a Stat is written as its icon and its
number in a rectangle --- \statplate{body}{3}, \statplate{mind}{2},
\statplate{face}{4}. Moving an object with your thoughts is still a
\keyword{Body} action because the object's weight resists you; scrying
on a secret is \keyword{Mind}, and swaying a guard's will is
\keyword{Face}. For more magical examples, see \emph{Which Base Stat?}
in \emph{Magic}.

\subsection*{Weird: When Reality Resists}
When nothing resists but reality itself, use \keyword{Weird} instead:
conjuring something from nothing or teleporting across a gap, rather
than crossing it. Weird equals your Magic rank minus one, from
\dice{0} to \dice{3}; you cannot raise it with Stat Points. Use Body,
Mind, or Face if matter, a puzzle, or a person's will is what pushes
back. See \emph{The Weird Stat} in \emph{Magic} for examples.

\subsection*{Stats from Objects}
% REAL RULE (Notion, 2026-09-21): an item can supply the base stat, using
% its tier. "Higher of the two" is the working idea, flagged for testing.
A base Stat does not have to come from your body. When you act
\emph{through} a thing --- a car, a rifle, a warhorse, a cutting deck ---
that object can supply the base Stat itself, using its \keyword{tier} in
place of your own number. Take the higher of the two. This changes the
dice you roll, not the size of your own consequence pools.

\begin{example}
Kira drives a rank-\rating{3} car through a chase. Her own \keyword{Body} is
\dice{2}, but the car's tier is \dice{3}, so the car supplies the base. With
\atrtag{skill}{Driving} and the first-Skill bonus, she rolls \dice{5} dice.
\end{example}

\section{Tags}
A \keyword{tag} is a word in an oval describing an edge you have on an
action. Count each tag that clearly applies and add one die for it. Tags
from Body, Stuff, Skills, and Magic can all count on the same roll. Their
small marks show where they came from. Both \atrtag{body}{Strong} and
\atrtag{magic}{Fire} add a die when they apply.

\begin{example}
Mara swings a sword she can barely lift. Her \atrtag{body}{Strong} and
\atrtag{skill}{Blades} both help: \keyword{Body} \dice{3} plus two tags
is \dice{5} dice, and Blades is her first Skill on the roll, so she adds
\dice{1} more for \dice{6} (see \emph{Your First Skill Adds +1}, below).
\end{example}

\subsection*{Skills and Magic Are Tags Too}
% REAL RULE, 2026-09-22: a Skill IS a tag --- not a separate resolution
% path, and not a ranked one either. You know it or you don't; a known
% skill adds one die, same as any other tag (plus the once-per-roll
% first-Skill bonus, 2026-09-23).
A \keyword{Skill} is a tag like any other on your sheet --- the trained
kind. You either know it or you don't; there is no individual rank to
raise. \atrtag{skill}{Athletics} and \atrtag{skill}{Blades} each add
one die when they apply, and the first Skill on a roll adds one extra
(see below). Skills cover everything mundane and
learned --- driving, acrobatics, lore, stealth, and combat.

\subsection*{Your First Skill Adds +1}
% REAL RULE, 2026-09-23 (clarified): every applicable Skill tag adds one
% die like any tag, AND the first applicable Skill adds one extra die on
% top. N applicable Skills = N+1 dice. The bonus counts toward the
% five-dice tag cap. Playtesters (02 03 05 06 07 09 10) misread the old
% wording as "only one Skill counts", so this spells out the arithmetic.
% Supersedes the earlier (never-implemented) idea of a $-1$ penalty for
% acting without a matching skill.
A Skill has to clearly fit the action, exactly like any other tag
(see \emph{Tags}, above), and the GM judges whether it does. Every
Skill that fits adds \dice{1} die, like any tag. On top of that, if
\emph{at least one} Skill fits, add \dice{1} more die, once per roll.
Your first Skill is worth two dice and each one after it is worth one:

\begin{center}
\begin{tabular}{@{}lccccc@{}}
Skills that fit & 0 & 1 & 2 & 3 & 4\\
Dice from Skills & 0 & 2 & 3 & 4 & 5\\
\end{tabular}
\end{center}

The bonus die counts toward the five-tag limit (see \emph{Five Tags,
Maximum}, below), so tags and the bonus together never add more than
\dice{5} dice. There is no penalty for having no applicable Skill at
all; you just miss out on the bonus.

\begin{example}
Mara has no \atrtag{skill}{Larceny}, but her \atrtag{skill}{Craft}
clearly covers picking a lock. The GM agrees it applies, so it adds
\dice{2} dice: one for the tag and one for being her first Skill. If her
\atrtag{skill}{Stealth} also applied because she was working silently
with a guard asleep next door, it would add \dice{1} more, for \dice{3}
dice from Skills.
\end{example}

% REAL RULE (Notion, 2026-09-21, decided): magic pool = base stat + one
% point per applicable tag. This SUPERSEDES the older "two dice per tag"
% rule.
Magic uses the same roll: choose the base Stat by what resists, then
count each applicable Magic tag. See \emph{Magic} for the tag lists and
what spells those tags can describe.

\begin{example}
Rooster has \keyword{Body} \dice{2} and the Magic tags
\atrtag{magic}{Destroy} and \atrtag{magic}{Water}. Shattering a window
with a hurled fist of water is a physical blow, so the base is
\keyword{Body}. Both tags apply: \dice{2}~+~\dice{2}~=~\dice{4} dice.
\end{example}

\subsection*{Five Tags, Maximum}
% REAL RULE, 2026-09-22: one flat cap for the whole system, replacing
% Magic's old per-rank "Tags per Spell" column. The first-Skill bonus
% counts toward this cap (decided 2026-09-23), so max pool stays 10.
Across all sources, a single roll never adds more than \dice{5} dice
from tags to its pool, and that includes the first-Skill bonus. This is
the one limit that applies everywhere --- Body, Stuff, Skills, and Magic
alike. If Mara lunges across a table with \atrtag{body}{Strong},
\atrtag{item}{Reach}, \atrtag{skill}{Blades}, and
\atrtag{skill}{Athletics} all applying, she has \dice{5} dice from tags
(four tags plus the bonus), which is already the most any roll can add.

\begin{notebox}[Stats and the Four Columns]
The four columns of character creation are not the same thing as the three
Stats. \keyword{Body}, \keyword{Skills}, and \keyword{Magic} ranks tell you
how much you may buy in each area; \keyword{Stuff} governs the gear you
carry. The Stats --- Body, Mind, Face --- are what you actually roll.
\end{notebox}

\section{Consequences}
% REAL RULE, 2026-09-22: each Stat has a companion pool (size = Stat
% rating). A shared party Coherence pool (starts at 5) absorbs overflow.
Failure is not just a missed roll --- it costs something. Each of your
three ordinary Stats has a companion pool the same size as its rating:
a \keyword{Body} pool, a \keyword{Mind} pool, and a \keyword{Face} pool.
Weird and object tiers do not give you extra pools or increase their
sizes. When a consequence costs you a point, take it from the pool the
fiction threatens. Usually it matches the Stat you rolled, but a
failed Mind roll to outwit an armed foe might cost you Body instead.

\begin{statblock}[Your Three Pools]
\begin{tabularx}{\linewidth}{@{}p{0.22\linewidth} Y@{}}
\keyword{Body Pool} & Physical harm --- wounds, exhaustion, injury.\\[3pt]
\keyword{Mind Pool} & Mental strain --- fear, confusion, a different kind of exhaustion.\\[3pt]
\keyword{Face Pool} & Social standing --- humiliation, broken trust, burned bridges.\\
\end{tabularx}
\end{statblock}

In a conflict, a successful action can cost an opponent a point from
one of these tracks, while a failure can cost you one. A partial success
can do both. The track affected need not match the Stat used to roll;
see \emph{Conflicts} in \emph{The Shape of a Session} for the full
exchange.

\subsection*{When a Pool Is Empty}
At \dice{0} in a pool, subtract \dice{1} die from every roll using its
Stat until you recover at least one point in that pool. At \dice{0}
Mind, for example, you roll one fewer die on Mind actions. This also
applies when an item's tier supplies the base for that Stat: a better
tool does not erase your exhaustion or shaken confidence. The penalty
does not stack if further consequences hit the empty pool; those losses
go to Coherence instead. Weird has no pool and no such penalty.

\subsection*{Recovering Your Pools}
When you enter a new world, fully refill your Body, Mind, and Face pools
to their Stat ratings. Within a world, recovery follows what happens in
the story. The GM decides whether an event restores part or all of a
pool, never more than its full size. A long rest might restore some or
all of Body and Mind, at the GM's discretion. Rest alone does not
restore Face: repair your standing by changing your social situation,
making amends, or earning recognition for a completed mission.

\begin{example}
Mara loses Face when the court laughs at her. After she completes a
quest, the king recognises her before that same court; the GM restores
her Face pool in full.
\end{example}

\subsection*{Coherence: The Party's Last Line}
When a pool would lose a point and it is already empty, the loss comes
out of \keyword{Coherence} instead --- one shared pool for the whole
party, starting at \dice{5}. Coherence isn't tied to any one character;
it is the party's grip on the world they are currently standing in, and
everyone draws from the same reserve.

\todo{When does the shared Coherence pool refill? Unlike personal
Body, Mind, and Face pools, its recovery is not yet specified.}

If Coherence ever reaches \dice{0}, something catastrophic happens, at
the GM's discretion --- the world itself has stopped holding together
around the party. One reliable option: they are torn loose from the
realm entirely and pulled into the next one, a clean and dramatic way to
close out a session. See \emph{For the Game Master} for more on running
this.
